\chapter{Source digests II: current approaches}
\label{ch:digests-modern}

This chapter continues Chapter~\ref{ch:digests-foundations}.  It digests the
survey of Dudko and the recent papers on satellite, Feigenbaum, parabolic,
neutral, hybrid, and decoration phenomena.  The conventions are the same: the
definitions and results are those of the source, the status is the source's own,
and the dependencies are stated explicitly.

\section{The survey's map of approaches}
\label{sec:digest-survey}

The survey~\cite{DudkoMLC} organizes the known approaches to MLC into three
broad regimes: the near-neutral regime, the quadratic-like regime, and the
puzzle regime, together with a near-degenerate regime.  Its listed approaches are
puzzles and parapuzzles, external rays and Böttcher coordinates, quadratic-like
renormalization, sector and cylinder renormalization in the near-parabolic and
near-Siegel cases, Molecule renormalization, a priori bounds, hybrid classes and
laminations, rigidity, combinatorial and kneading methods, pseudo-Siegel disks,
and virtual Julia sets.

\subsection*{Definitions}

\begin{definition}[Renormalization operator]
\label{def:survey-renorm}
For $p_c(z)=z^2+c$ and a small $\eta>0$, the renormalization is
\[
 [\mathcal R_\eta p_c](z)=\frac1\eta\bigl(p_c^{\,q}(\eta z+c)-c\bigr),
\]
where $q$ is the first return time to the $\eta$-neighbourhood of $c$.
\end{definition}

\begin{definition}[Satellite and primitive components]
\label{def:survey-satellite}
A hyperbolic component $\mathcal H$ is \emph{satellite} if its root has multiplier $e^{2\pi i p/q}$ with $q\ge2$,
so that $\mathcal H$ is attached to a larger component, and \emph{primitive} if its root has multiplier $1$.  The survey
states that satellite components have zero entropy and primitive components positive entropy (Figure~3 of the survey,
\texttt{refs/2512.24171v1.txt}, lines 179--182).
\end{definition}

\begin{definition}[Renorm-expanding and renorm-balanced regimes]
\label{def:survey-expanding}
In the renorm-expanding regime the derivatives of all globalized branches tend to
infinity.  In the renorm-balanced regime some branches converge geometrically to
a polynomial or transcendental map $G=\lim_n\mathcal R^{(q_n)}_{\eta_n}p_c$.
The regime is pure if the property holds for every sequence $\eta_n\to0$.
\end{definition}

\begin{definition}[Quadratic-like maps and hybrid classes]
\label{def:survey-ql}
A quadratic-like map is a degree-two branched covering $f\colon U\to V$ with
$U\Subset V$.  Its filled Julia set is $K_f=\{z\in U:\ f^n(z)\in U\ \forall n\ge1\}$,
and $J_f=\partial K_f$.  If $J_f$ is connected, $f$ is hybrid equivalent to a
unique $z^2+c$ with $c\in\M$.
\end{definition}

\begin{definition}[Rotation-number sets and neutral maps]
\label{def:survey-rotation}
Put $\Theta=\{[0;a_1,a_2,\dots]:a_i\ge1\}$, $\Theta_{\rm bnd}=\{[0;a_i]:a_i\le M_\theta\}$,
and $\Theta_{>N}=\{[0;a_i]:a_i>N\}$; $\overline\Theta$ allows $a_i=\infty$.
A neutral map is $f_\theta(z)=e^{2\pi i\theta}z+z^2$.
\end{definition}

\begin{definition}[Sector, cylinder, puzzle, and Molecule renormalizations]
\label{def:survey-other-renorm}
Sector renormalization is the first return to the smallest sector bounded by a
curve $\gamma$ and its image, with the sides identified dynamically.  Cylinder
renormalization is the quotient of a Lavaurs map by a parabolic map.  Puzzle
renormalization is the first return to a puzzle piece $Y$, restricted to the
critical piece $Y_1$ of degree two.  The survey unifies satellite and neutral
renormalizations as Molecule renormalization.
\end{definition}

\subsection*{Results}

\begin{theorem}[Yoccoz's reduction, as presented in the survey]
\label{thm:survey-yoccoz-reduction}
For a parameter $c$, if nested strictly decreasing small copies $\M_n$ are chosen
with $\bigcap_n\M_n=\{c\}$, then the survey reduces local connectivity at $c$ to a
rigidity statement for these copies.  The survey gives this as a reduction of MLC
to rigidity, not as a complete proof.
\end{theorem}

\begin{theorem}[Lyubich; codimension-one stability in the quadratic-like class]
\label{thm:survey-lyubich-codim}
There is an analytic function $\tau$ on a Banach neighbourhood such that $f$ and
$g$ are hybrid conjugate if and only if $\tau(f)=\tau(g)$.  Status: cited.
\end{theorem}

\begin{theorem}[Quadratic-like hyperbolicity]
\label{thm:survey-ql-hyp}
For every bound $\bar p$ on renormalization periods there is a hyperbolic
renormalization horseshoe with one-dimensional unstable manifolds, which attracts
exponentially all infinitely renormalizable quadratic-like maps with periods
$p_n\le\bar p$.  Status: stated in the survey as a theorem of its quadratic-like part.
\end{theorem}

\begin{theorem}[Near-parabolic hyperbolicity; Inou--Shishikura]
\label{thm:survey-near-parabolic}
Cylinder renormalization is hyperbolic for high-type $\theta$ with bounded
partial quotients, for $N$ large.  This implies sectorial bounds.  Status: cited
to Inou--Shishikura.  The survey does not define the set $\Theta_{\le N}$ explicitly,
so the precise class is not recorded here.
\end{theorem}

\begin{theorem}[Pacman and Siegel hyperbolicity]
\label{thm:survey-pacman}
For a bounded-type periodic rotation number $\theta$, Pacman/Siegel renormalization
is hyperbolic with one-dimensional unstable manifold.  The survey cites this to
Dudko, Lyubich and Selinger~\cite{DudkoLyubichSelinger}, and records that the golden
mean case was previously obtained with computer assistance by Gaidashev and Yampolsky.
\end{theorem}

\begin{theorem}[Bounded-type quadratic-like MLC; Dudko--Lyubich]
\label{thm:survey-bounded-ql-mlc}
For all bounded-type quadratic-like combinatorics, a priori bounds and MLC hold.
In particular MLC holds at the period-doubling Feigenbaum parameter.  Combined with
puzzle arguments the survey obtains MLC along a real line and along every
Mandelbrot vein.  Status: proved in~\cite{DudkoLyubichFeigenbaum}; the veins
statement is the survey's own consequence.
\end{theorem}

\begin{theorem}[Neutral a priori bounds; Dudko--Lyubich]
\label{thm:survey-neutral-bounds}
Uniform pseudo-Siegel a priori bounds hold for every neutral quadratic polynomial,
and the nested pseudo-Siegel disks converge to the Mother Hedgehog.  Status:
stated in the survey as a Dudko--Lyubich theorem.
\end{theorem}

\begin{theorem}[Anti-Molecule near-degenerate case]
\label{thm:survey-anti-molecule}
Positive-entropy near-degenerate estimates, combined with puzzle estimates, give a
priori bounds and MLC under the anti-Molecule condition.  Status: the survey's own statement, with the condition defined in its near-degenerate part.
\end{theorem}

\subsection*{Open problems listed by the survey}

\begin{openestimate}[Survey problems]
\label{prob:survey}
The survey lists the following as open: uniform bounded geometry of all hyperbolic
components; uniform a priori bounds for all $c\in\partial\M$ and all scales; a
single uniformly hyperbolic renormalization operator covering the quadratic-like,
near-neutral, and puzzle regimes; hyperbolicity of neutral renormalization for all
$\overline\Theta$; hyperbolicity of Molecule renormalization; pseudo-Siegel bounds
in the remaining unbounded satellite quadratic-like cases; and a virtual Molecule
version of the near-degenerate regime.  The survey states that the last two
problems, together with the other arguments it lists, should imply full MLC.
\end{openestimate}

\subsection*{Logical relations stated in the survey}

\begin{itemize}
\item Bounded geometry of hyperbolic components implies MLC, by shrinking nested
  copies and Yoccoz's reduction.
\item Quadratic-like hyperbolicity follows from bounded-type MLC and a priori bounds
  together with Lyubich's quadratic-like theory.
\item The neutral a priori bounds imply a compact sectorial renormalization operator.
\item Positive-entropy estimates and puzzle estimates give MLC in the anti-Molecule case.
\item The Feigenbaum theorem together with puzzle arguments gives MLC along real lines
  and Mandelbrot veins.
\end{itemize}

\section{Dudko--Lyubich: satellite parameters}
\label{sec:digest-satellite}

The satellite paper~\cite{DudkoLyubichSatellite} proves that $\M$ and the Julia sets $J_{c_n}$ are locally connected at a
constructed sequence of infinitely renormalizable parameters of bounded satellite type, and that $J_{c_n}$ has positive area.
The full digest, with the definitions and the scope, is in Section~\ref{sec:dl-satellite}.

\section{Dudko--Lyubich: MLC at Feigenbaum points}
\label{sec:digest-feigenbaum}

The paper~\cite{DudkoLyubichFeigenbaum} proves beau bounds at every Feigenbaum parameter, that is, at every infinitely
renormalizable parameter of bounded type, and deduces MLC there from Lyubich's rigidity theorem.  The full digest is in
Section~\ref{sec:dl-feigenbaum}.

\section{Kahn--Kapiamba--Lyubich: parabolic families}
\label{sec:digest-kkl}

The paper~\cite{KahnKapiambaLyubich} proves beau bounds and super-optimism for parabolically bounded families of prime
primitive types, and MLC pointwise on these families under one input from an unpublished manuscript.  The full digest is
in Section~\ref{sec:kkl}.

\section{Neutral and hybrid sources}
\label{sec:digest-neutral-hybrid}

\subsection*{Dudko--Luo--Lyubich: $\psi^\bullet$-quasi-Siegel maps}

\begin{definition}[Pseudo-Siegel bounds]
\label{def:pseudo-siegel}
For $f_\theta(z)=e^{2\pi i\theta}z+z^2$, the pseudo-Siegel bounds are nested disks
$\widehat Z^{-1}\supseteq\widehat Z^0\supseteq\cdots\supseteq H\ni0$, where $H$ is
the filled postcritical set, $f_\theta$ is injective on each $\widehat Z^m$, and
$\bigcap_m\widehat Z^m=H$.  The constants are uniform in $\theta$.
\end{definition}

\begin{definition}[$\psi^\bullet$-quasi-Siegel map]
\label{def:psi-ql-siegel}
A $\psi^\bullet$-quasi-Siegel map is a pair $F=(f,\iota)$ of holomorphic maps between
disks with the properties stated in~\cite{DudkoLuoLyubichNeutral}, including that
$f$ is a double branched cover with a unique critical point on the boundary of its
domain of definition.
\end{definition}

\begin{theorem}[Dudko--Luo--Lyubich, Theorem~1.1 ($\psi^\bullet$-case)]
\label{thm:dll-psi}
If $F$ is a $\psi^\bullet$-quasi-Siegel map with an eventually golden-mean Siegel
disk and degeneration $K_F$, there is a nested sequence of geodesic pseudo-Siegel
disks satisfying the pseudo-Siegel bounds.  The outermost disk is a quasidisk with
constant depending on $K_F$.  Status: proved in the paper.  Its regularization result is conditional on the
combinatorial hypotheses stated there and is imported from an earlier
Dudko--Lyubich paper, which the source cites.
\end{theorem}

\subsection*{Dudko--Lim: pseudo-bubbles}

\begin{definition}[Bubbles and pseudo-bubbles]
\label{def:bubbles}
The generation-zero bubble is $B^0_\theta=\overline{f_\theta^{-1}(Z_\theta)\setminus Z_\theta}$,
and its pseudo-bubble is the closure of the component of $f_\theta^{-1}(\operatorname{int}\widehat Z_\theta)$
disjoint from $Z_\theta$.  A generation-$g$ bubble is a component of $f_\theta^{-g}(B^0_\theta)$,
with the corresponding pseudo-bubble.  Its root is the unique critical point of
$f_\theta^{g+1}$ in it.
\end{definition}

\begin{theorem}[Dudko--Lim, pseudo-bubble bounds]
\label{thm:dl-bubbles}
There is a universal $K>1$ such that every pseudo-bubble is a $K$-quasidisk.  If the
root is $c_{-k}$ with $q_n\le k<q_{n+1}$, then
\[
 K^{-1}|c_{-k}-c_{q_n-k}|\le\operatorname{diam}(B)\le\operatorname{diam}(\widehat B)
 \le K|c_{-k}-c_{q_n-k}|,
\]
and every point of $\widehat B$ lies within hyperbolic distance $K$ of the geodesic
landing at $c_{-k}$.  Status: proved in the paper from the earlier pseudo-Siegel bounds of Dudko and
Lyubich (cited in the source as~[DL22]) and standard Koebe and quasiconformal
estimates.  The paper states no MLC theorem.
\end{theorem}

\subsection*{Dudko--Lim--Lyubich: rigidity of the neutral attractor}

\begin{definition}[Neutral cascade and butterfly]
\label{def:neutral-cascade}
A neutral cascade is the limiting transcendental semigroup of a bi-infinite
renormalization tower; its parameter semigroup is commutative and totally ordered
by time.  A butterfly $\mathcal B_n$ is a pair of maps $F^{Q_n+Q_{n-1}}\colon\mathcal U'_n\to\mathcal W_n$
and $F^{Q_n}\colon\mathcal U''_n\to\mathcal W_n$, with depth-$m$ enrichments and
$\mathcal B^\infty_n=\bigcup_{m\ge n}\mathcal B^m_n$.
\end{definition}

\begin{theorem}[Conformal rigidity of bi-infinite towers]
\label{thm:dll-conformal}
Combinatorially equivalent bi-infinite towers in the renormalization attractor are
conformally conjugate on definite neighbourhoods of their Mother Hedgehogs.  Status:
proved in the paper.
\end{theorem}

\begin{theorem}[Affine rigidity of neutral cascades]
\label{thm:dll-affine}
The quasiconformal conjugacy in Theorem~\ref{thm:dll-conformal} is affine.  Status:
proved in the paper, using a dynamical-hairiness theorem and quasiconformal
Thurston equivalence.
\end{theorem}

\subsection*{Shen: hybrid classes of generalized polynomial-like maps}

\begin{definition}[Generalized polynomial-like maps, hybrid classes, criticality]
\label{def:shen}
A generalized polynomial-like map $f\colon U\to V$ has finite unions of Jordan disks
$U\Subset V$ as domain and range, each component proper onto a component of $V$, and
analytic extension across $\overline U$ with no critical point on $\partial U$.  Its
filled Julia set is $K(f)=\{z\in U:\ f^n(z)\in U\ \forall n\ge0\}$.  The criticality of $f$ is
$d(f)=\sum_{c\in\operatorname{Crit}(f)\cap K(f)}\operatorname{ord}(f,c)$.
\end{definition}

\begin{theorem}[Shen, Theorem~1.1]
\label{thm:shen}
For $f_0$ in the relevant class, the hybrid class $\mathcal H_*(f,\mathcal V)$ is,
locally and among maps with constant criticality, an embedded complex analytic
submanifold of the Banach space of the class, closed in $\mathcal V$, of codimension
$d(f_0)$.  Status: proved in the paper.  The paper's proof uses a polynomial
result of Astorg, subordination, and a generalized straightening theorem of
Baranski.  It notes that earlier proofs for generalized polynomial-like maps relied
on approximation by hyperbolic maps, an unresolved issue.
\end{theorem}

\section{Dudko: the decoration theorem}
\label{sec:digest-decoration}

The paper~\cite{DudkoDecoration} studies the components of $\M$ removed by a small
copy.

\begin{definition}[Small copy and decorations]
\label{def:decoration}
A quadratic map is Douady--Hubbard renormalizable of period $p$ if there is a
dividing repelling cycle with associated domains $Y^p_1\Subset Y^0_1$, and $c$ does not
escape from $Y^p_1$ under $f^p$.  The small copy $\M_s$ is the closure of the
parameters satisfying this condition with a fixed ray pattern.  The component of
$\M\setminus\M_s$ containing the main cardioid is big; the other components are
decorations.
\end{definition}

\begin{theorem}[Decoration theorem]
\label{thm:decoration}
For every $\varepsilon>0$ there are at most finitely many connected components of
$\M\setminus\M_s$ of diameter at least $\varepsilon$.  Status: proved in the paper
by puzzle and parapuzzle methods.  The paper also states the same result for the
multibrot sets $\M_d=\{c:\ J(z^d+c)\text{ connected}\}$, $d\ge2$.  It does not
state global MLC.
\end{theorem}

\section{Kahn--Lyubich: Elephant Eyes}
\label{sec:digest-elephant}

The paper~\cite{KahnLyubichElephant} proves estimates for elephant-eye combinatorics, with $q$ unbounded and $p=q+O(1)$,
and states local connectivity and MLC for these parameters through external inputs.  The full digest, including the status of
its unnumbered main statement, is in Section~\ref{sec:kahn-IV}.

\section{Dependency map}
\label{sec:digests-dependencies}

Table~\ref{tab:digest-dependencies} collects the dependencies stated above.  An
entry names an input, a result that uses it, and the status of that result.

\begin{table}[ht]
\centering\small
\begin{tabularx}{\textwidth}{@{}>{\raggedright\arraybackslash}p{0.23\textwidth}>{\raggedright\arraybackslash}X>{\raggedright\arraybackslash}p{0.27\textwidth}@{}}
\textbf{Result} & \textbf{Inputs used} & \textbf{Status in the source}\\
\hline
Yoccoz finite renormalization; boundaries of hyperbolic components &
Puzzle and shrinking arguments &
cited\\
Lyubich I--III (Section~\ref{sec:digest-lyubich-1993}) &
straightening, Yoccoz, Sullivan, Hu--Jiang--Jiang, McMullen &
proved\\
Survey, quadratic-like codimension one and hyperbolicity &
Lyubich's quadratic-like theory &
cited or stated\\
Pacman/Siegel hyperbolicity &
bounded-type periodic rotation numbers; Gaidashev--Yampolsky for the golden mean &
Dudko--Lyubich--Selinger; see Chapter~\ref{ch:digests-pacman-neutral}\\
Bounded-type quadratic-like MLC &
a priori bounds; puzzle arguments &
proved in~\cite{DudkoLyubichFeigenbaum}\\
Satellite bounded-type MLC &
molecule surgery near periodic irrational rotations &
proved for the paper's class\\
Feigenbaum Main A and B &
beau bounds; bounded relative periods &
proved\\
KKL Main A, B, C &
beau bounds; puzzle Teichm\"uller diameter; \emph{Elephant puzzles} &
proved, with B depending on an unpublished input\\
Pseudo-Siegel bounds (DLL) &
regularization from an earlier paper &
proved, conditional on stated combinatorics\\
Pseudo-bubble bounds (DL) &
pseudo-Siegel bounds; Koebe and quasiconformal estimates &
proved\\
Conformal and affine rigidity (DLL) &
pseudo-Siegel bounds; dynamical hairiness; Thurston equivalence &
proved\\
Hybrid classes (Shen) &
Astorg; subordination; Baranski straightening &
proved\\
Decoration theorem &
puzzles and parapuzzles; Yoccoz's trivial-fiber result &
proved\\
Elephant Eyes a priori bounds &
thin--thick decomposition; amplification; pull-off &
proved\\
Elephant Eyes MLC &
Hu--Jiang, Jiang; Lyubich; Ecalle--Voronin cylinders &
the paper's own statement\\
\end{tabularx}
\caption{Dependencies among the digested results.}
\label{tab:digest-dependencies}
\end{table}

\section{What these digests do not yet cover}
\label{sec:digests-gaps}

These digests are not a complete review of the literature.  The following gaps
are known and recorded here, so that they are not mistaken for coverage.
\begin{itemize}
\item The definitions of quasi-hyperbolic parameters, of the Douady--Hubbard
  renormalization of polynomials, and of external rays are not restated in this
  chapter.  They belong to the foundations chapter and to~\cite{DH}.
\item Digested in Chapter~\ref{ch:digests-apriori-limbs}: the Kahn--Lyubich series on a priori bounds (Parts~I--IV), the
  Dudko--Lyubich satellite and Feigenbaum papers, the Kahn--Kapiamba--Lyubich paper, and the Dudko--Schleicher paper on limbs.
\item Several inputs are cited but not digested: Sullivan's complex bounds, Hu--Jiang
  and Jiang, McMullen, Inou--Shishikura, Gaidashev--Yampolsky, and the unpublished
  \emph{Elephant puzzles}.  Their statements are used only as recorded by the
  sources that cite them.
\item Digested in Chapters~\ref{ch:digests-monographs}, \ref{ch:digests-pacman-neutral}, and
  \ref{ch:digests-laminations-entropy}: Milnor's lectures, Lyubich's ICTP notes, the Pacman paper of Dudko,
  Lyubich and Selinger, the neutral a priori bounds of Dudko and Lyubich, the matings-with-laminations and
  core-entropy papers, the Sierpi\'nski-carpet paper of Dudko and Luo (Chapter~\ref{ch:digests-sierpinski}), and Słodkowski's holomorphic-motion theorems.  The following are not yet digested:
  Lyubich's online volumes I--II (a web search found no freely available Volumes~III--IV of
  \emph{Conformal Geometry and Dynamics of Quadratic Polynomials}), the papers \cite{LyubichDynamics,LyubichIII}; the Dudko papers on
  canonical decompositions and correspondences, which are
  indirectly related to MLC and are listed in the reference inventory; and the Bartholdi--Dudko series on
  algorithmic aspects of branched coverings, which concerns computability rather than MLC.
\item The survey's notation $\Theta_{\le N}$ is not defined explicitly in the survey, so
  the class in its near-parabolic statement is recorded only as the survey states it.
\item Lean formalization of these sources is not complete.  The repository contains
  interfaces for several of the hypotheses above, as described in
  Chapter~\ref{ch:conditional-matrix} and Appendix~\ref{ch:formal-map}.  It contains
  no formal proof of Main A or Main B of any of the recent papers, no quasiconformal
  surgery, and no straightening theorem.  The exact examples of Milnor's lectures
  (Chapter~\ref{ch:digests-monographs}) and the real slice are formalized.
\end{itemize}
